\section{Full-source pointwise reconstruction and the arithmetic remainder}
\label{sec:v43-full-source-inversion}
We use the correction of Section~\ref{sec:v43-complete-correction}
without assuming that its full mixed residual transform is integrable.
The order of inversion matters: take the discrete label coefficient
first, and then invert its roof transform. Each roof inverse is
absolutely convergent. This gives a reconstruction of the full law,
including the source left outside the guarded packet.

\subsection{Labelwise inversion with a prescribed pointwise error}
Retain $e_\ell,q_\ell,b_\ell$ and $\varepsilon$ from
Theorem~\ref{thm:v43-complete-correction}. Define the exact component
transform and its residual by
\begin{equation}\label{eq:v43-component-transform}
 \Phi_\ell^w(b)=\int_{Y_R^*}\one_{\{(K_n,N_n)=(k,m)\}}
                      w(x)e^{ibT_n(x)}\,d\nu_R^*(x),\qquad
 h_\ell(b)=\Phi_\ell^w(b)-\widehat e_\ell(b)=\widehat q_\ell(b).
\end{equation}
The physical record has not been convolved. Its finite component
transform is given by the same exact finite collision graph as its
germs. Put $A_{j,\ell}=\|\partial_t^j q_\ell\|_1$ for $j=0,2$.
The individual $A_{2,\ell}$ are finite; their unweighted sum is not
assumed finite.

\begin{theorem}[Complete raw inversion in coefficient-first order]
\label{thm:v43-labelwise-inversion}
For every component of positive reference mass,
\begin{equation}\label{eq:v43-component-fourier-norm}
 \int_\R|h_\ell(b)|\,db\le4\sqrt{A_{0,\ell}A_{2,\ell}},\qquad
 f_\ell(t)=e_\ell(t)+\frac1{2\pi}\int_\R e^{-ibt}h_\ell(b)\,db.
\end{equation}
The equality holds almost everywhere, at every regular density point,
and with the right trace at every finite-trace edge. At such an edge
the point value of $e_\ell$ is assigned its right trace, while
$q_\ell$ uses its continuous representative. At a divergent germ it is an equality of densities off the singular
value, not an assertion of a finite density there.

For $\delta>0$ choose, when $b_\ell>0$,
\begin{equation}\label{eq:v43-label-band}
 \Lambda_\ell=\max\left\{1,\frac{A_{2,\ell}}{\pi\delta b_\ell}\right\},
 \qquad
 p_\delta^w(\ell,t)=e_\ell(t)
       +\frac1{2\pi}\int_{-\Lambda_\ell}^{\Lambda_\ell}
                                            e^{-ibt}h_\ell(b)\,db.
\end{equation}
For a null component put $p_\delta^w=0$. Then
\begin{equation}\label{eq:v43-summed-pointwise-error}
 \sum_\ell\operatorname*{ess\,sup}_{t\in\R}
             |f_\ell(t)-p_\delta^w(\ell,t)|\le\delta.
\end{equation}
For a fixed finite label set one may use the maximum of its
$\Lambda_\ell$ as a single bandwidth. The bounds do not omit grazing,
section decisions, normal states, or trajectories with any prescribed
actual collision label.
\end{theorem}
\begin{proof}
The $W^{2,1}$ residual satisfies
\[
 |h_\ell(b)|\le\min\{A_{0,\ell},A_{2,\ell}|b|^{-2}\}.
\]
The integral of the right side is $4\sqrt{A_{0,\ell}A_{2,\ell}}$.
If either norm is zero, the compactly supported residual is zero.
Absolute Fourier inversion gives its continuous representative. The
zero value and derivative traces in Lemma~\ref{lem:v43-small-jet-mass}
make the trace conventions in the statement consistent with the
intrinsic germs. Integration over $|b|>\Lambda$ gives the exact bound
$A_{2,\ell}/(\pi\Lambda)$. Substitution of
\eqref{eq:v43-label-band} bounds it by $\delta b_\ell$, and summing
$\sum b_\ell=1$ proves \eqref{eq:v43-summed-pointwise-error}.
Increasing a bandwidth preserves this upper bound.

There is no unjustified interchange of two infinite integrals.
By Corollary~\ref{cor:v43-correction-transform}, the full residual
transform has the uniformly absolutely convergent discrete expansion
\[
 \widehat Q^\varepsilon(u,s,b)
     =\sum_{k,m}e^{i(u\cdot k+sm)}h_{(k,m)}(b).
\]
For each fixed $b$, its torus coefficient is therefore exactly
$h_\ell(b)$. Only that coefficient is subsequently integrated in
$b$, where \eqref{eq:v43-component-fourier-norm} applies. Joint
$L^1(\T^3\times\R)$ integrability of the transform is neither
used nor obtained by this argument.
\end{proof}

\begin{corollary}[A prescribed reconstruction scale at long return times]
\label{cor:v43-prescribed-reconstruction}
Fix $P>0$. Choose
$\varepsilon_n=\delta_n=(1+n)^{-P-3}$ in the preceding construction.
Then
\begin{equation}\label{eq:v43-reconstruction-rate}
 \sup_{R\in I} n^2\sum_\ell\operatorname*{ess\,sup}_{t}
        |p_{n,R}^w(\ell,t)-p_{\delta_n}^w(\ell,t)|
                  \le n^2(1+n)^{-P-3}.
\end{equation}
Here the preparation data and the bandwidths are allowed to depend
on $n,R,\ell,w$. For central return targets it suffices to retain
labels $m\le L_n$, where
\begin{equation}\label{eq:v43-explicit-count-cutoff}
 L_n=\left\lceil\lambda n+\frac{P+4}{c_0}\log(2+n)\right\rceil,
 \qquad \lambda>\max\{c^{-1}+1,a/c_0+1\}.
\end{equation}
On these retained labels the reconstruction error contains no
collision-tail density. The total variation of the omitted correction
and residual is bounded by \eqref{eq:v43-correction-count-tail}.
\end{corollary}
\begin{proof}
Equation~\eqref{eq:v43-summed-pointwise-error} has constant one at
every $R$, proving the first assertion with the indicated
parameter-dependent choices. A fixed return central window has
$m=n/c+O(\sqrt n)$; its labels are at most $L_n$ eventually.
Restrictions by the collision label are exact, so a discarded label
cannot contribute to the density of a retained label. Finally
$an-c_0L_n\le-c_0n-(P+4)\log(2+n)$, up to the harmless strengthening
from the strict choice of $\lambda$. Apply the preceding tail bound.
\end{proof}

This is a prescribed error for \emph{reconstruction}, not a prescribed
rate for a Gaussian local limit. In particular the formula for
$\Lambda_\ell$ retains the actual germ derivative budget; it provides
no polynomial or exponential upper bound for that budget as $n$
grows. The global mass budget for $E^{\varepsilon_n}$ can be made
smaller than $n^{-2}$ while its pointwise height remains large.
These statements must not be interchanged.

\subsection{The arithmetic main term for the unchanged singleton}
We now take $w=1$. Fix once a real even function
$\chi\in C_c^\infty((-1,1))$, with $0\le\chi\le1$ and
$\chi=1$ near zero. For a \emph{fixed} $B>0$ put
\[
 K_B(t)=\frac1{2\pi}\int e^{-ibt}\chi(b/B)\,db,
 \qquad \int K_B=1.
\]
The kernel is real and even, but need not be nonnegative. The
band-limited test theorem applies to it without a positivity
assumption. Define the full-source signed correction
\begin{align}\label{eq:v43-fixed-band-correction}
 \mathscr D_{B,n,R}(\ell,t)
  ={}&e_\ell(t)-(K_B*e_\ell)(t)\notag\\
     &+\frac1{2\pi}\int_\R e^{-ibt}
                         (1-\chi(b/B))h_\ell(b)\,db.
\end{align}
This integral is absolute componentwise. All measures in this formula
refer to the full actual law at the stated return index.

\begin{theorem}[Arithmetic raw closure with an evaluated central term]
\label{thm:v43-arithmetic-raw-criterion}
At every fixed $B>0$ there is the exact almost-everywhere identity
\begin{equation}\label{eq:v43-correction-identity}
 \mathscr D_{B,n,R}(\ell,t)
       =p_{n,R}(\ell,t)-(K_B*p_{n,R}(\ell,\cdot))(t).
\end{equation}
In particular the correction is independent of the germ localization
radii and of $\varepsilon$. The same identity applies to all finite
truncations once the specified label is retained.

Write $y=(k_1,k_2,t,n)$ and
$Z=(y-m(0,0,\bar\tau_R,c))/\sqrt m$, in that order. Uniformly in
$R$, $m\ge n\ge1$, and $|Z|\le M$,
\begin{equation}\label{eq:v43-arithmetic-raw-decomposition}
 m^2p_{n,R}(k,m,t)
   =\mathcal L_{m,R}(y)+m^2\mathscr D_{B,n,R}(k,m,t)+o(1).
\end{equation}
The error is in essential supremum in the roof target and ordinary
supremum in the other displayed parameters. The transition kernel is
exactly the evaluated finite Gaussian peak sum of
Theorem~\ref{thm:v41-uniform-transition}, with its arithmetic branches
retained. Consequently the radius-uniform pointwise transition local
law holds if and only if
\begin{equation}\label{eq:v43-single-remainder-criterion}
 \sup_{R,n,k}\ m^2\operatorname*{ess\,sup}_{t:\ |Z|\le M}
                          |\mathscr D_{B,n,R}(k,m,t)|\longrightarrow0
\end{equation}
for every $M$. If it holds for one fixed $B>0$, it holds for every
fixed $B>0$.

At a fixed radius, the corresponding main term is
$c\mathfrak a_R(k,n,m)g_{\Omega_R}(Z)$; equivalently in the return
normalization it is
$\mathfrak a_R(k,n,m)g_{D_R}(V_n)$. No zero-residue assertion or
replacement of the prescribed singleton by a packet is needed for
these statements.
\end{theorem}
\begin{proof}
The absolutely convergent component inverse gives
\[
 \frac1{2\pi}\int e^{-ibt}(1-\chi(b/B))h_\ell(b)\,db
                    =q_\ell(t)-(K_B*q_\ell)(t).
\]
Add the edge difference and use $e_\ell+q_\ell=p_{n,R}(\ell,\cdot)$.
This proves \eqref{eq:v43-correction-identity}. If a different
extraction is used, it gives the identical right side, so the
invariance is exact, not an asymptotic comparison of absolute bounds.

Because $K_B$ is even, its convolution at $t$ is precisely
$P^{K_B}_{n,m,R}(k,t)$ with the sign convention
$K_B(T_n-t)$. This is one fixed integrable test, with integrable
Fourier transform supported in $[-B,B]$. Equation
\eqref{eq:v41-uniform-test} and its vanishing imaginary part give
\[
 m^2(K_B*p_{n,R}(\ell,\cdot))(t)=\mathcal L_{m,R}(y)+o(1)
\]
uniformly even without restricting the targets to a central compact
set. Substitute \eqref{eq:v43-correction-identity}. This proves
\eqref{eq:v43-arithmetic-raw-decomposition} and both implications in
\eqref{eq:v43-single-remainder-criterion}. Comparing the identity at
two fixed bands proves the last band-independence assertion. Each
band is fixed before $m\to\infty$.

At fixed radius use the specialization of
Theorem~\ref{thm:v41-uniform-transition} and the covariance change in
\eqref{eq:v40-arithmetic-original-normalization}. On the central
sets $n/m=c+O(m^{-1/2})$, so the same conversion applies to the error
and to the signed correction. The arithmetic coefficient is uniformly
bounded by the finite resonance order. When it is zero, the exact
component measure is zero by
Theorem~\ref{thm:v40-arithmetic-return-LLT}; it is not assigned a
conditional density law with a zero denominator.
\end{proof}

\begin{corollary}[A finite integral certificate for the remaining error]
\label{cor:v43-finite-remainder-certificate}
Replace the last integral of \eqref{eq:v43-fixed-band-correction} by
its restriction to $|b|\le\Lambda_\ell$, where
\[
 \Lambda_\ell\ge\max\{B,1,A_{2,\ell}/(\pi\delta b_\ell)\}.
\]
Call the resulting quantity $\mathscr D^{\rm fin}_{B,n,R}$. Then
\[
 \sum_\ell\operatorname*{ess\,sup}_t
       |\mathscr D_{B,n,R}-\mathscr D^{\rm fin}_{B,n,R}|(\ell,t)
                         \le\delta.
\]
With $\delta_n=(1+n)^{-P-3}$ this difference is negligible at the
local scale on every central set. Thus the necessary and sufficient
criterion may be tested on this finite, full-source inverse with an
explicit tail certificate. No spectral power estimate is evaluated
at the varying $\Lambda_\ell$ in this assertion.
\end{corollary}
\begin{proof}
Beyond $\Lambda_\ell\ge B$, the factor $1-\chi(b/B)$ is exactly
one. The omitted integral is at most
$A_{2,\ell}/(\pi\Lambda_\ell)\le\delta b_\ell$. Sum the estimate.
For central targets $m/n$ is bounded, so $m^2\delta_n\to0$.
The finite integral is left as an exact physical quantity; its smallness
is not deduced from the fixed-band spectral theorem.
\end{proof}

\paragraph{The remaining analytical estimate.}
The complete correction now exists, with all boundary germs included,
in a topology with an explicit summable mass budget. Its full-source
pointwise reconstruction has a certified error after the label is
fixed. What is not yet estimated is the signed finite inverse in
Corollary~\ref{cor:v43-finite-remainder-certificate} on the central
long-time scale. The old guarded derivative bound cannot be inserted
as a polynomial bound for the new germ budgets, and the small edge
mass cannot be substituted for its pointwise height. A proof of
\eqref{eq:v43-single-remainder-criterion}, including the arithmetic
transition kernel rather than an assumed unit residue, remains the
precise raw-density closure step. The definition of that step is now
independent of count cutoffs, auxiliary source guards and choices of
singular localization.
